\section{Introduction}
\label{sec:intro}

Redistricting ensemble methods have emerged as a principled approach to
evaluating and generating congressional maps. By sampling broad regions of the
legally feasible plan space, methods such as ReCom \citep{deford2021recombination}
and GerryChain \citep{duchin2019gerrychain} provide a reference distribution
against which any proposed map can be measured for partisan or racial bias.
The core primitive is a \emph{recombination step}: two adjacent districts are
merged and re-split along a random spanning tree, producing a new feasible plan
in $O(n)$ time per step.

\paragraph{The bipartition failure problem.}
This approach works well when the number of districts $k$ is small or when
$k$ decomposes into small factors. But many states have seat counts that are
large primes or contain large prime factors: Texas has $k = 38 = 2 \times 19$,
Pennsylvania has $k = 17$ (prime), New York has $k = 26 = 2 \times 13$,
and Michigan has $k = 13$ (prime). A direct $k$-way ReCom ensemble on such
states requires, at each step, merging two adjacent districts and re-splitting
the merged subgraph into a balanced bipartition. The correct failure mode is
geometric, not a balance-ratio issue: in a $k=38$ state, merging two adjacent
districts produces a region whose spanning tree must achieve a $1:1$ balance
\emph{within} the merged region, but the surrounding $k-2 = 36$ districts
constrain the boundary so severely that balanced spanning tree cuts become rare.
GerryChain \emph{without pair reselection} stalls in this regime; GerryChain
\emph{with pair reselection} can succeed by trying different adjacent pairs
until a balanced cut is found, at the cost of lower throughput in Python.

\paragraph{BisectionEnsemble: local 2-way ReCom at every node.}
We propose a structural fix: instead of running a $k$-way ensemble at the top
level, embed 2-way ensemble sampling \emph{inside} the bisection tree. The
bisection tree (see ApportionRegions \citep{dellaLibera2026apportion} and
GeoSection \citep{dellaLibera2026geosection}) already decomposes the redistricting
problem into a hierarchy of 2-way cuts. At each node of this tree---which
manages a subregion of roughly $n/2^\ell$ tracts at depth $\ell$ and must split it into
2 balanced parts---we replace the single METIS call with a $T$-step local
ReCom chain seeded by the METIS bisection.

Because every sampled ensemble node is 2-way, the ReCom sampler is never asked
to solve a large-arity split such as a 19-way ApportionRegions split. It only
attempts a $1:1$ balance within the current node's subregion (or near-$1:1$ for
floor/ceil splits). This avoids the large-arity failure mode; it does not prove
that every finite local chain will accept quickly, so the algorithm logs
rejections and falls back to the METIS initialisation if no local proposal is
accepted. METIS calls at $p'>2$ ApportionRegions nodes are outside this
guarantee. The bisection tree structure is fully preserved, retaining population
balance and contiguity through the same mechanisms as the original bisection
algorithm.

\paragraph{Key advantages.}
\begin{enumerate}
  \item \textbf{Large-arity split avoidance at bisection nodes}: the ReCom
        sampler is used only for 2-way node splits, while non-binary
        ApportionRegions nodes remain METIS responsibilities.
  \item \textbf{Local feasibility sampling}: Each node samples the local
        feasible space of bisections of its subregion. At depth $\ell$, each
        node manages $\approx n/2^\ell$ tracts; the locality advantage grows
        with depth. The root node ($\ell=0$) manages the full state, while
        leaf nodes manage $\approx n/k$ tracts.
  \item \textbf{Parallelisable}: Tree nodes at the same depth are independent;
        each node's ensemble can run on a separate CPU core. For a balanced
        binary tree of depth $d = \log_2 k$, level $\ell$ has $2^\ell$ nodes
        running in parallel.
  \item \textbf{Structure-preserving}: The bisection tree guarantees population
        balance and contiguity at every level, regardless of which bisection cut
        is selected by the ensemble percentile.
\end{enumerate}

\paragraph{Percentile selection and legal posture.}
The percentile parameter $p \in [0, 1]$ selects which bisection cut to use from
the ensemble distribution. At $p = 0.5$ (median), we select the plan at the
median edge-cut rank---the most ``typical'' feasible bisection of the subregion.
At $p = 0.0$ (minimum), we select the minimum edge-cut plan---the most compact
bisection. The legal argument for percentile selection is identical to that of
PercentileSweep \citep{dellaLibera2026percentilesweep}: the choice of $p$ is
a transparent, pre-committed procedure that does not depend on partisan outcomes.

\paragraph{Contributions.}
\begin{enumerate}
  \item \textbf{Algorithm}: BisectionEnsemble---a node-local 2-way ReCom
        ensemble embedded in the bisection tree, with percentile plan selection.
  \item \textbf{Implementation}: \texttt{SeedCompositor::BisectionEnsemble\{p, ensemble\_steps\}}
        in \texttt{redist-cli}, using \texttt{split\_subgraph\_bisection\_ensemble()}.
  \item \textbf{Failure-mode scoping}: proof that BisectionEnsemble avoids the
        large-arity full-plan split inside the ReCom sampler, with explicit
        finite-$T$ fallback behavior.
  \item \textbf{Empirical evaluation}: Comparison with standard METIS bisection
        and GerryChain on NC ($k=14$), WI ($k=8$), and TX ($k=38$).
\end{enumerate}
